% front_abbreviations
\clearpage
\phantomsection
\section*{\bgen{СПИСЪК НА СЪКРАЩЕНИЯТА}{LIST OF ABBREVIATIONS}}
\addcontentsline{toc}{section}{\bgen{СПИСЪК НА СЪКРАЩЕНИЯТА}{LIST OF ABBREVIATIONS}}

Списъкът съдържа съкращенията, използвани в текста, и само тях. Имената на системни
извиквания, на функции и на типове не са съкращения и не се превеждат; те се дават с
шрифт за код на мястото на употребата им.

\begin{longtable}{@{}L{3.0cm}L{11.5cm}@{}}
\toprule
\textbf{Съкращение} & \textbf{Значение} \\
\midrule
\endfirsthead
\toprule
\textbf{Съкращение} & \textbf{Значение} \\
\midrule
\endhead
\bottomrule
\endfoot

ALPN & Application-Layer Protocol Negotiation, договаряне на протокола на приложно ниво в рамките на TLS ръкостискането \cite{rfc7301} \\
API & Application Programming Interface, програмен интерфейс \\
BCa & Bias-corrected and accelerated, вид бутстрап доверителен интервал \cite{efron1987better} \\
eBPF & Extended Berkeley Packet Filter, механизъм за изпълнение на проверени програми в ядрото на Linux \\
HPACK & Компресия на заглавни полета в HTTP/2 \cite{rfc7541} \\
HTML & HyperText Markup Language \\
HTTP & HyperText Transfer Protocol \\
HTTPS & HTTP над TLS \\
IOCP & Input/Output Completion Ports, механизъм за вход и изход по завършване в Windows \\
IP & Internet Protocol \\
JSON & JavaScript Object Notation \\
MTU & Maximum Transmission Unit, най-голям размер на предаваната единица \\
QPACK & Компресия на заглавни полета в HTTP/3 \cite{rfc9114} \\
QUIC & Транспортен протокол над UDP с вградена криптография \cite{rfc9000} \\
RFC & Request for Comments, поредица от документи на IETF \\
TCP & Transmission Control Protocol \\
TLS & Transport Layer Security \\
UDP & User Datagram Protocol \\
XSS & Cross-Site Scripting, вид уязвимост при вграждане на неекранирани данни \\
\end{longtable}

\vspace{0.5em}
\noindent
Термините, оставени на английски в текста по правилата за терминология, не са съкращения и
затова не влизат тук. Такива са \emph{handler}, \emph{middleware}, \emph{callback},
\emph{stream}, \emph{backend} и \emph{endpoint}.
